\section{Auswertung}
\label{sec:Auswertung}

Zunächst werden in der Vorbereitungsaufgabe für verschiedene Zylinderlängen 
(Reihen aus jeweils $50 \unit{mm}$- und $75 \unit{mm}$-Zylindern)
Frequenzspektren aufgenommen. Beispielhafte Spektren sind in \autoref{fig:1}
abgebildet.

\begin{figure}[H]
    \centering
    \begin{subfigure}[b]{0.45\textwidth}
        \centering
        \includegraphics[width=\textwidth]{build/2vor1.pdf}
        \caption{Spektrum $50 \unit{mm}$}
        \label{fig:sub1}
    \end{subfigure}
    \hfill
    \begin{subfigure}[b]{0.45\textwidth}
        \centering
        \includegraphics[width=\textwidth]{build/2vor6.pdf}
        \caption{Spektrum $100 \unit{mm}$}
        \label{fig:sub2}
    \end{subfigure}
    \vspace{1em}
    
    \begin{subfigure}[b]{0.45\textwidth}
        \centering
        \includegraphics[width=\textwidth]{build/2vor12.pdf}
        \caption{Spektrum $150 \unit{mm}$(aus $50 \unit{mm}$ Zylindern)}
        \label{fig:sub3}
    \end{subfigure}
    \hfill
    \begin{subfigure}[b]{0.45\textwidth}
        \centering
        \includegraphics[width=\textwidth]{build/vor1.pdf}
        \caption{Spektrum $75 \unit{mm}$}
        \label{fig:sub4}
    \end{subfigure}
    \vspace{1em}
    
    \begin{subfigure}[b]{0.45\textwidth}
        \centering
        \includegraphics[width=\textwidth]{build/vor4.pdf}
        \caption{Spektrum $150 \unit{mm}$(aus $75 \unit{mm}$ Zylindern)}
        \label{fig:sub5}
    \end{subfigure}
    \hfill
    \begin{subfigure}[b]{0.45\textwidth}
        \centering
        \includegraphics[width=\textwidth]{build/vor8.pdf}
        \caption{Spektrum $225 \unit{mm}$}
        \label{fig:sub6}
    \end{subfigure}
    \caption{Beispielhafte Spektren für verschiedene Zylinderlängen.}
    \label{fig:1}
\end{figure}
\noindent Aus zwei dieser Spektren, in diesem Fall die des $50 \unit{mm}$- und des $75 \unit{mm}$-Zylinders,
wird die Schallgeschwindigkeit bestimmt. Das geschieht auf Grundlage der Interferenzbedingung stehender Wellen 
in einem Zylinder. Stehende Wellen bilden sich in einem Zylinder bei konstruktiver Interferenz 
unter folgender Bedingung aus:
\begin{equation}
    L = n \cdot \frac{\lambda}{2} \quad \text{mit } n \in \mathbb{N}
    \label{eq:interf}
\end{equation}
\noindent Dabei ist $\lambda$ die Wellenlänge und $L$ die Zylinderlänge. Umgestellt nach $\lambda$ ergibt sich die folgende Gleichung:
\begin{equation}
    \lambda = \frac{L\cdot 2}{n}
    \label{eq:lambda}
\end{equation}
Des Weiteren gilt der Zusammenhang zwischen Frequenz $f$ und $\lambda$:
\begin{equation}
    f = \frac{c}{\lambda}
    \label{eq:freq}
\end{equation}
Dabei ist $c$ die Schallgeschwindigkeit.
\noindent Daraus lässt sich über die gemittelten Abstände unserer Peaks in den
besagten zwei Spektren des $50 \unit{mm}$- und $75 \unit{mm}$-Zylinders die
Schallgeschwindigkeit berechnen. Durch Umstellen von \autoref{eq:freq} nach 
$\lambda$ und anschließendes einsetzen in \autoref{eq:lambda} ergibt sich 
folgender Ausdruck:
\begin{equation}
    \label{eq:1}
    f_\text{n} = n \cdot \frac{c}{2\cdot L}
\end{equation}
So kann ein Ausdruck der Proportionalität zwischen den Abständen der Peaks und
der Schallgeschwindigkeit (\autoref{eq:12}) gebildet werden:
\begin{equation}
    \Delta f = f_{n+1} - f_n = (n+1) \cdot \frac{c}{2L} - n \cdot \frac{c}{2L} = \frac{c}{2L}
    \label{eq:12}
\end{equation}
Die Frequenzen, welche positive Interferenz aufzeigen, werden mit der Funktion \texttt{find\_peaks} 
aus der Bibliothek SciPy \cite{virtanen2020scipy} bestimmt und dann deren Abstände gemittelt. 
Die Maxima sind in \autoref{fig:13} markiert.
\begin{figure}[H]
    \centering
    \includegraphics[width=\textwidth]{build/peaks.pdf}
    \caption{frequenzspektrum des $50$, sowie $75 \unit{mm}$-Zylinders, sowie gefundene Maxima}
    \label{fig:13}
\end{figure}
\noindent Für den $50 \unit{mm}$-Zylinder ergeben sich folgende Abstände und daraus
resultierende Schallgeschwindigkeiten (\autoref{eq:werte1}):
\begin{align}
    \label{eq:werte1}
    \overline{f_\text{n}-f_\text{n+1}} &= \qty{3.27(0.23)e3}{\hertz}\\
    c_\text{50mm} &= \qty{327.10(22.75)}{m/s}
\end{align}
Für den $75 \unit{mm}$-Zylinder ergeben sich die Werte in \autoref{eqn:werte}:
\begin{align}
    \label{eqn:werte}
    \overline{f_\text{n}-f_\text{n+1}} &= \qty{2.19(0.17)e3}{\hertz}\\
    c_\text{75mm} &= \qty{329.10(24.87)}{m/s}
\end{align}

\subsection{Darstellung des Wasserstoffatoms}
\label{sec:02}
In diesem Abschnitt wird die Abhängigkeit der Amplitude der Druckwelle vom Polarwinkel 
im Kugelresonator untersucht und mit der quantenmechanischen Beschreibung des Wasserstoffatoms verglichen. Das geschieht über 
die Kugelflächenfunktionen, welche die Polar- und Azimutalabhängigkeit beim Wasserstoffatom beschreiben. 
\vspace{0.5em} \\
\noindent Zunächst wird ein hochaufgelöstes Spektrum bei einem eingestellten Winkel von $180^\circ$ aufgenommen, um die 
Resonanzfrequenzen zu identifizieren. Bei der Wasserstoffanalogie wäre dies das 
Abtasten nach den radialen Eigenfrequenzen, also den Werten, die die Quantenzahl $n$ annehmen kann. 
Der radiale Anteil der Schrödingergleichung ist jedoch nicht gleich dem radialen Anteil der Helmholtzgleichung,
also lässt sich mit diesem Spektrum keine Aussage über die Lage der Energieeigenwerte des Wasserstoffatoms machen. 
Es dient lediglich dazu, radiale Eigenfrequenzen zu finden, welche dann in ihrer polaren bzw. azimutalen Abhängigkeit weiter untersucht werden können.
Das Übersichtsspektrum ist in \autoref{fig:uebersicht} abgebildet.
\begin{figure}[H]
    \centering
    \includegraphics[width=\textwidth]{build/spektrum.pdf}
    \caption{Übersichtsspektrum bei Winkel von $180^\circ$}
    \label{fig:uebersicht}
\end{figure}
\noindent Die gefundenen Eigenfrequenzen bei $f = \qty{2.3}{\kilo\hertz}, \qty{3.7}{\kilo\hertz}, \qty{5.0}{\kilo\hertz}$ und $\qty{7.4}{\kilo\hertz}$ werden 
in $5^\circ$-Schritten von einem Winkel von $0^\circ$ bis $180^\circ$ untersucht. Dazu wird für jeden dieser Winkel ein Spektrum 
von $\qty{0.4}{\kilo\hertz}$ um die jeweilige Eigenfrequenz herum aufgenommen. Die Amplituden der Peaks aus diesen Spektren werden wieder 
mit der Funktion \texttt{find\_peaks} \cite{virtanen2020scipy} ausgewertet. In \autoref{fig:kff} sind die Amplituden gegen den Polarwinkel für die jeweilige 
Eigenfrequenz dargestellt.
Es ist zu beachten, dass die Achse des Lautsprechers und des Mikrofons in einem
$45^\circ$-Winkel zu der Ebene steht, in der der Winkel $\alpha$ verändert 
wird. Der gemessene Winkel muss mittels \autoref{eqn:transf} in den Polarwinkel transformiert werden. 
\begin{equation}
    \label{eqn:transf}
    \theta = \arccos\left(\frac{1}{2} \cos\alpha - \frac{1}{2}\right)
\end{equation}

\begin{figure}[H]
    \centering
    \begin{subfigure}[b]{0.45\textwidth}
        \centering
        \includegraphics[width=\textwidth]{build/legendre_11_fit.png}
        \caption{Polarplot für $f = \qty{2.3}{\kilo\hertz}$}
        \label{fig:kff_1}
    \end{subfigure}
    \hfill
    \begin{subfigure}[b]{0.45\textwidth}
        \centering
        \includegraphics[width=\textwidth]{build/legendre_12_fit.png}
        \caption{Polarplot für $f = \qty{3.7}{\kilo\hertz}$}
        \label{fig:kff_2}
    \end{subfigure}
    \vspace{1em}
    
    \begin{subfigure}[b]{0.45\textwidth}
        \centering
        \includegraphics[width=\textwidth]{build/legendre_13_fit.png}
        \caption{Polarplot für $f = \qty{5.0}{\kilo\hertz}$}
        \label{fig:kff_3}
    \end{subfigure}
    \hfill
    \begin{subfigure}[b]{0.45\textwidth}
        \centering
        \includegraphics[width=\textwidth]{build/legendre_14_fit.png}
        \caption{Polarplot für $f = \qty{7.4}{\kilo\hertz}$}
        \label{fig:kff_4}
    \end{subfigure}
    \caption{Amplituden der untersuchten Resonanzfrequenzen in Abhängigkeit vom Polarwinkel (Vergleich mit Kugelflächenfunktionen).}
    \label{fig:kff}
\end{figure}


\subsection{Aufspaltung der $l$-Entartung}
Bisher wurde für jedes $n$, also jede Energiequantenzahl bzw. jede radiale Eigenfrequenz 
in dem Kugelresonator, nur das Orbital mit der jeweils größten Drehimpulsquantenzahl 
$l$ in den Polarplots identifiziert. Das lag daran, dass sich die entarteten Zustände der jeweiligen radialen 
Eigenfrequenz in einer Überlagerung befinden. In dieser Überlagerung dominieren die 
Zustände mit den größten Drehimpulsquantenzahlen, weil diese noch mehr entartete $m$-Zustände aufweisen.
Es werden in den Kugelresonator Zwischenringe mit den folgenden Dicken eingebaut:
\begin{equation}
    d_1 = \qty{3}{\milli\meter}, \quad d_2 = \qty{6}{\milli\meter} \quad \text{und zusammen} \quad d_3 = \qty{9}{\milli\meter}.
\end{equation}
\noindent Durch diese wird die Kugelsymmetrie gebrochen und es entsteht eine Aufspaltung der 
radialen Eigenfrequenzen in die jeweiligen Zustände der Drehimpulsquantenzahlen für dieses $n$. 
Des Weiteren ist zu beachten, dass durch die Symmetriebrechung die Ebene, in der die Winkel abgetastet werden, nun 
genau der Polarebene entspricht, wir also direkt den Polarwinkel messen.
Es werden wie in \autoref{sec:02} für Winkel zwischen $0^\circ$ und $180^\circ$ die Amplituden 
der Druckwelle im Bereich von $\qty{0.6}{\kilo\hertz}$ um die jeweilige radiale Eigenfrequenz herum aufgenommen.
Diese Amplituden werden dann für die zwei verschiedenen $l$ in \autoref{fig:auf1} und \autoref{fig:auf2} dargestellt.
\begin{figure}[H]
    \centering
    \begin{subfigure}[b]{0.45\textwidth}
        \centering
        \includegraphics[width=\textwidth]{build/aufspaltung_l1_m0.png}
        \caption{Polarplot für $l = 0$}
        \label{fig:auf1}
    \end{subfigure}
    \hfill
    \begin{subfigure}[b]{0.45\textwidth}
        \centering
        \includegraphics[width=\textwidth]{build/aufspaltung_l1_m1.png}
        \caption{Polarplot für $l = 1$}
        \label{fig:auf2}
    \end{subfigure}
    \caption{Aufspaltung der radialen Eigenfrequenzen in Abhängigkeit der Drehimpulsquantenzahl $l$.}
    \label{fig:aufspaltung_main}
\end{figure}
